% Options for packages loaded elsewhere
\PassOptionsToPackage{unicode}{hyperref}
\PassOptionsToPackage{hyphens}{url}
\PassOptionsToPackage{dvipsnames,svgnames*,x11names*}{xcolor}
%
\documentclass[
  11pt]{article}
\usepackage{lmodern}
\usepackage{amssymb,amsmath}
\usepackage{ifxetex,ifluatex}
\ifnum 0\ifxetex 1\fi\ifluatex 1\fi=0 % if pdftex
  \usepackage[T1]{fontenc}
  \usepackage[utf8]{inputenc}
  \usepackage{textcomp} % provide euro and other symbols
\else % if luatex or xetex
  \usepackage{unicode-math}
  \defaultfontfeatures{Scale=MatchLowercase}
  \defaultfontfeatures[\rmfamily]{Ligatures=TeX,Scale=1}
\fi
% Use upquote if available, for straight quotes in verbatim environments
\IfFileExists{upquote.sty}{\usepackage{upquote}}{}
\IfFileExists{microtype.sty}{% use microtype if available
  \usepackage[]{microtype}
  \UseMicrotypeSet[protrusion]{basicmath} % disable protrusion for tt fonts
}{}
\makeatletter
\@ifundefined{KOMAClassName}{% if non-KOMA class
  \IfFileExists{parskip.sty}{%
    \usepackage{parskip}
  }{% else
    \setlength{\parindent}{0pt}
    \setlength{\parskip}{6pt plus 2pt minus 1pt}}
}{% if KOMA class
  \KOMAoptions{parskip=half}}
\makeatother
\usepackage{xcolor}
\IfFileExists{xurl.sty}{\usepackage{xurl}}{} % add URL line breaks if available
\IfFileExists{bookmark.sty}{\usepackage{bookmark}}{\usepackage{hyperref}}
\hypersetup{
  pdftitle={A Quantum World Model for Model-Based Reinforcement Learning},
  pdfauthor={V. Liutarevich},
  colorlinks=true,
  linkcolor=NavyBlue,
  filecolor=Maroon,
  citecolor=Blue,
  urlcolor=NavyBlue,
  pdfcreator={LaTeX via pandoc}}
\urlstyle{same} % disable monospaced font for URLs
\usepackage[margin=1in]{geometry}
\usepackage{longtable,booktabs}
% Correct order of tables after \paragraph or \subparagraph
\usepackage{etoolbox}
\makeatletter
\patchcmd\longtable{\par}{\if@noskipsec\mbox{}\fi\par}{}{}
\makeatother
% Allow footnotes in longtable head/foot
\IfFileExists{footnotehyper.sty}{\usepackage{footnotehyper}}{\usepackage{footnote}}
\makesavenoteenv{longtable}
\setlength{\emergencystretch}{3em} % prevent overfull lines
\providecommand{\tightlist}{%
  \setlength{\itemsep}{0pt}\setlength{\parskip}{0pt}}
\setcounter{secnumdepth}{-\maxdimen} % remove section numbering
\usepackage{booktabs}
\usepackage{longtable}
\usepackage{amsmath}
\usepackage[font=small,labelfont=bf]{caption}
\usepackage{float}
\setlength{\emergencystretch}{3em}
\renewcommand{\arraystretch}{1.15}
\usepackage{fancyhdr}
\pagestyle{fancy}
\fancyhf{}
\fancyfoot[C]{\thepage}
\fancyhead[L]{\footnotesize\itshape Draft v1 — 2026-10-07.}
\renewcommand{\headrulewidth}{0.2pt}

\title{A Quantum World Model for Model-Based Reinforcement Learning}
\author{V. Liutarevich}
\date{}

\begin{document}
\maketitle
\begin{abstract}
Quantum reinforcement learning has so far placed parameterised quantum
circuits in the policy or the value function. We place one in the
\textbf{world model}, which we argue is the better-motivated target: a
world model predicts dynamics, and dynamics is what quantum systems
natively express.

We present, to our knowledge, the first integration of a parameterised
quantum circuit into the dynamics model of a model-based reinforcement
learning agent, trained end-to-end. The circuit replaces the prior of a
recurrent state-space model --- the component that predicts the next
stochastic latent from the recurrent state. Because latent imagination
rolls the model forward on the prior alone, with no access to
observations, \textbf{the policy is trained entirely on trajectories
generated by a quantum circuit.}

We describe what the integration requires and release an implementation
together with a verification harness. The resulting agent learns: on
CartPole it reaches an evaluation return of 318.6 against a
random-policy return of approximately 22, with policy improvement driven
entirely by imagined rollouts from the quantum prior. \textbf{{[}RUN ---
a budget-matched comparison against the capacity-matched classical
control is outstanding; current runs differ in training budget and
neither arm reached the solved threshold. No performance claim is made
until it is complete.{]}}

We further characterise the design space through ablations over qubit
count, circuit depth, entangler topology, readout locality, and an
entanglement-removed control \textbf{{[}RUN{]}}, and discuss what the
architecture enables.
\end{abstract}

\hypertarget{introduction}{%
\subsection{1. Introduction}\label{introduction}}

Reinforcement learning with quantum components has, to date, been almost
entirely model-free. Parameterised quantum circuits have been used as
policies and as value-function approximators, inheriting the sample
inefficiency of the model-free setting while adding the cost of circuit
evaluation.

Model-based reinforcement learning offers a different and, we argue,
better insertion point. A model-based agent learns a \emph{world model}
--- a predictive model of environment dynamics --- and then optimises
its policy against that model rather than against the environment. The
world model is the component whose job is to represent dynamics, which
is precisely what a quantum system does naturally, and it is the
component where a quantum substrate has the clearest motivation.

This paper places a parameterised quantum circuit inside the world model
of a model-based agent and trains the whole system end-to-end.
Specifically, the circuit replaces the \textbf{prior} of a recurrent
state-space model (RSSM): the distribution over the next stochastic
latent given the deterministic recurrent state. This choice is not
incidental. In latent imagination the agent rolls the model forward
using the prior alone, without observations, and trains its policy on
the resulting trajectories. Replacing the prior therefore means the
agent's imagined experience --- all of the data its policy ever sees
during optimisation --- is generated by a quantum circuit.

\textbf{Contributions.}

\begin{enumerate}
\def\labelenumi{\arabic{enumi}.}
\tightlist
\item
  \textbf{A quantum dynamics prior for recurrent state-space models},
  and the first model-based reinforcement learning agent with a quantum
  world model trained end-to-end (Section 4).
\item
  \textbf{An account of what the integration requires}: mapping circuit
  readouts to categorical latent logits, maintaining gradient flow
  through multi-step imagination, and the interaction with KL balancing
  (Sections 4--5).
\item
  \textbf{A verification harness} of integration and correctness checks,
  released with the implementation, addressing failure modes that are
  silent rather than loud (Section 5).
\item
  \textbf{Evidence of viability}: the agent learns through imagination
  from the quantum prior, reaching an evaluation return of 318.6 on
  CartPole against a random-policy baseline of approximately 22.
\item
  \textbf{A characterisation of the design space} through ablations
  \textbf{{[}RUN{]}}.
\end{enumerate}

We do not claim a performance advantage. The contribution is that the
architecture exists, trains, and works; our comparison against the
classical prior is reported with intervals and is consistent with
parity. For a first system measured against a mature and heavily tuned
baseline, we regard parity as the expected and appropriate outcome.

\begin{center}\rule{0.5\linewidth}{0.5pt}\end{center}

\hypertarget{background-recurrent-state-space-models-and-latent-imagination}{%
\subsection{2. Background: recurrent state-space models and latent
imagination}\label{background-recurrent-state-space-models-and-latent-imagination}}

A model-based agent of the DreamerV3 family \textbf{{[}1{]}} maintains a
latent state factorised into a deterministic recurrent component \(h_t\)
and a stochastic component \(z_t\) represented as a set of categorical
variables. Three distributions define the world model:

\begin{itemize}
\tightlist
\item
  a \textbf{recurrent update} \(h_t = f(h_{t-1}, z_{t-1}, a_{t-1})\),
\item
  a \textbf{posterior} \(q(z_t \mid h_t, o_t)\), which incorporates the
  observation,
\item
  a \textbf{prior} \(p(z_t \mid h_t)\), which does not.
\end{itemize}

The world model is trained to reconstruct observations and rewards while
keeping the prior close to the posterior under a balanced KL term, so
that the prior learns to predict what the posterior would have inferred
had an observation been available.

\textbf{Imagination.} Policy optimisation does not interact with the
environment. The agent samples a starting latent from replayed
experience and rolls the world model forward for a fixed horizon using
\textbf{the prior only}, since no future observations exist. The actor
and critic are trained on these imagined trajectories via
\(\lambda\)-returns.

The prior is therefore the generative core of the agent: it is what the
agent dreams with. This makes it the natural and most consequential
place to introduce a quantum component, and it is the choice we make.

\begin{center}\rule{0.5\linewidth}{0.5pt}\end{center}

\hypertarget{related-work}{%
\subsection{3. Related work}\label{related-work}}

\textbf{Quantum reinforcement learning.} Existing work places
parameterised quantum circuits in the policy or value function ---
variational quantum circuits as deep Q-network approximators, as policy
networks for discrete and continuous control, and related model-free
constructions \textbf{{[}2, 3{]}}. To our knowledge no prior work places
a quantum circuit in the dynamics model of a model-based agent.
\textbf{{[}AUDIT --- this is the paper's central novelty claim and
requires a systematic search before submission, not a sample of it.
Record the search strings, dates, and databases.{]}}

\textbf{World models for quantum problems.} The converse direction
exists: model-based reinforcement learning has been applied \emph{to}
quantum problems, for instance DreamQAS \textbf{{[}4{]}}, which learns a
classical world model over circuit-construction trajectories for quantum
architecture search. That work uses a classical model to design quantum
circuits; we use a quantum circuit inside a model.

\textbf{Quantum sequence models.} Quantum attention and transformer
architectures \textbf{{[}5, 6{]}} provide the token-mixing construction
our dynamics prior adapts. Recent evaluations \textbf{{[}7, 8{]}} find
that quantum models at simulable scale generally match rather than
exceed classical baselines of comparable capacity, which informs the
expectations we set in Section 8.

\textbf{Model-based reinforcement learning.} DreamerV3 \textbf{{[}1{]}}
and the RSSM line it belongs to provide the architecture we extend.

\begin{center}\rule{0.5\linewidth}{0.5pt}\end{center}

\hypertarget{architecture}{%
\subsection{4. Architecture}\label{architecture}}

\hypertarget{the-quantum-dynamics-prior}{%
\subsubsection{4.1 The quantum dynamics
prior}\label{the-quantum-dynamics-prior}}

In the baseline RSSM the prior is a multilayer perceptron mapping the
deterministic recurrent state to logits over the categorical stochastic
latent:

\[ \text{logits} = \mathrm{MLP}(h_t) \in \mathbb{R}^{S \times C}, \qquad
   p(z_t \mid h_t) = \mathrm{Categorical}(\text{logits}) \]

for \(S\) stochastic variables with \(C\) classes each. We replace this
map with a quantum circuit, preserving the interface exactly so that
every other component of the agent --- recurrent update, posterior,
decoders, actor, critic, KL balancing, imagination --- is unchanged.

The quantum prior comprises three stages:

\begin{enumerate}
\def\labelenumi{\arabic{enumi}.}
\tightlist
\item
  \textbf{Down-projection.} A linear map
  \(h_t \in \mathbb{R}^{256} \mapsto \mathbb{R}^{6}\) followed by
  \(\tanh\) and scaling by \(\pi\), producing rotation angles in
  \([-\pi, \pi]\) for \(n = 6\) data qubits.
\item
  \textbf{The circuit.} Angle encoding with data re-uploading over
  \(L = 4\) layers, brickwork entangler, local Pauli-\(Z\) readout,
  giving 6 outputs in \([-1, 1]\). Trainable circuit parameters number
  \(L \times (n+1) = 28\).
\item
  \textbf{Expansion head.} A two-layer map from the 6 readouts to the
  \(S \times C = 256\) categorical logits, through a hidden width of 64
  with a SiLU nonlinearity.
\end{enumerate}

\hypertarget{the-bottleneck-and-an-honest-accounting}{%
\subsubsection{4.1.1 The bottleneck, and an honest
accounting}\label{the-bottleneck-and-an-honest-accounting}}

The expansion head is necessary --- the latent requires 256 logits and
the circuit emits 6 --- but it is not small. The parameter breakdown of
the substituted module is:

\begin{longtable}[]{@{}ll@{}}
\toprule
component & parameters\tabularnewline
\midrule
\endhead
down-projection (\(256 \rightarrow 6\)) & 1,542\tabularnewline
\textbf{quantum circuit} & \textbf{28}\tabularnewline
expansion head (\(6 \rightarrow 64 \rightarrow 256\)) &
17,088\tabularnewline
\textbf{total (quantum prior)} & \textbf{18,658}\tabularnewline
classical control, same positions & 18,672\tabularnewline
original DreamerV3 prior (replaced) & 131,840\tabularnewline
\bottomrule
\end{longtable}

The circuit is 0.15\% of the module's parameters. We state this plainly
because it bounds what the architecture can be claimed to demonstrate,
and because the natural reading --- that a quantum circuit is doing the
world-modelling --- requires qualification.

The defensible claim is one of \textbf{information bottleneck rather
than parameter share}: every prediction the prior makes, and therefore
every imagined transition the policy is trained on, passes through six
real numbers produced by a quantum circuit. The head cannot route around
it. Whether the circuit or the head carries the predictive work is
nonetheless an empirical question, and we answer it by varying the head
width (Section 7).

\hypertarget{why-the-prior-and-not-elsewhere}{%
\subsubsection{4.2 Why the prior, and not
elsewhere}\label{why-the-prior-and-not-elsewhere}}

Three positions in the agent could host a quantum circuit: the policy,
the value function, or the world model. The first two are what the
existing literature has used. We take the world model for three reasons.

\textbf{It is the component that models dynamics.} The motivation for a
quantum substrate rests on the claim that quantum systems are natural
representations of dynamical processes. A policy is a map from states to
actions; it carries no such motivation.

\textbf{It is where errors compound, and so where model quality matters
most.} The agent's policy is trained entirely on imagined rollouts.
Improvements or degradations in the prior propagate through the horizon
and are amplified, making the prior the component with the greatest
leverage over final performance.

\textbf{It is the smallest intervention with the largest reach.}
Replacing the prior changes \(L \times (n+1)\) parameters
\textbf{{[}VERIFY --- state the count{]}} out of the agent's total, yet
every imagined trajectory the policy ever sees passes through it.

\hypertarget{the-classical-control}{%
\subsubsection{4.3 The classical control}\label{the-classical-control}}

To isolate the effect of the quantum substrate from the architectural
change accompanying it, we compare against a classical prior occupying
the same position with the same shape: a \(256 \rightarrow 6\)
bottleneck, a \(6 \rightarrow 6\) map in place of the circuit, and the
identical \(6 \rightarrow 64 \rightarrow 256\) expansion head. This
follows the capacity-matched control protocol of \textbf{{[}5{]}}, and
it matches to 14 parameters out of 18,672 --- the two arms differ in the
substituted map and nothing else.

We note that this control is deliberately \emph{not} DreamerV3's
original prior network, which has 131,840 parameters. Both of our arms
are therefore substantially narrower than the baseline architecture at
that position, which is a property of the comparison to keep in view
when reading Section 6.

\begin{center}\rule{0.5\linewidth}{0.5pt}\end{center}

\hypertarget{implementation-and-verification}{%
\subsection{5. Implementation and
verification}\label{implementation-and-verification}}

\hypertarget{what-the-integration-requires}{%
\subsubsection{5.1 What the integration
requires}\label{what-the-integration-requires}}

Three aspects proved non-obvious and are worth recording.

\textbf{Optimiser reconstruction.} The quantum prior is installed by
substituting the prior module within an already-constructed agent. An
optimiser built before the substitution holds references to the
\emph{replaced} parameters; the new circuit parameters are silently
untracked, and training proceeds without error while the quantum
component never updates. The optimiser must be rebuilt after
substitution. This failure is silent, produces plausible learning
curves, and is not detectable from loss values alone.

\textbf{Gradient flow through imagination.} The actor is trained on
trajectories produced by rolling the prior forward for \(H\) steps.
Gradients therefore reach the circuit through a composition of \(H\)
applications of the same circuit, and verifying that they arrive
non-zero is a necessary check rather than a formality.

\textbf{KL balancing.} The world model loss regularises the prior toward
the posterior with asymmetric weighting. A prior with a materially
different parameterisation changes the effective strength of this term,
and we report the resulting KL trajectories \textbf{{[}RESULTS{]}}.

\hypertarget{verification-harness}{%
\subsubsection{5.2 Verification harness}\label{verification-harness}}

We release a suite of integration and physics checks that must pass
before any training run. It verifies, among others, that circuit
parameters receive non-zero gradients, that the substituted prior
produces distributions of correct shape and normalisation, that the
optimiser tracks the new parameters, that imagination rollouts remain
finite over the full horizon, and that quantum and classical arms have
matched parameter counts. \textbf{{[}VERIFY --- enumerate the full check
list from \texttt{selftest.py}.{]}}

The motivation is that the failure modes in this integration are silent.
A missing gradient path, an untracked parameter group, or a mis-shaped
distribution produce runs that complete normally and report plausible
numbers.

\begin{center}\rule{0.5\linewidth}{0.5pt}\end{center}

\hypertarget{experiments}{%
\subsection{6. Experiments}\label{experiments}}

\hypertarget{setup}{%
\subsubsection{6.1 Setup}\label{setup}}

CartPole, full observation, with a reduced training configuration
relative to DreamerV3's defaults: batch 8 against 16, sequence length 16
against 32, imagination horizon 10 against 15, and a gradient step every
4 environment steps against every 2. The reduction is forced by the cost
of circuit evaluation --- imagination performs \texttt{horizon}
sequential circuit calls over
\texttt{batch\ \$\textbackslash{}times\$\ seq\_len} states per gradient
step --- and applies identically to both arms. The solved threshold is a
mean evaluation return of 475.

\hypertarget{current-status-of-the-comparison}{%
\subsubsection{6.2 Current status of the
comparison}\label{current-status-of-the-comparison}}

\textbf{The runs completed to date do not support a performance claim,
and we report them here as a status rather than a result.}

\begin{longtable}[]{@{}llll@{}}
\toprule
arm & seed & env steps & final eval return\tabularnewline
\midrule
\endhead
classical control & 0 & 80,000 & 442.7 ± 101.4\tabularnewline
classical control & 1 & 40,000 & 186.1 ± 67.7\tabularnewline
classical control & 2 & 40,000 & 153.0 ± 46.2\tabularnewline
quantum prior & 0 & 80,000 & 318.6 ± 86.4\tabularnewline
quantum prior & 1 & \textbf{2,500} & 85.1 ± 27.4\tabularnewline
quantum prior & 2 & 40,000 & 145.7 ± 61.5\tabularnewline
\bottomrule
\end{longtable}

Two problems are evident. \textbf{The training budgets are not matched
across seeds} --- they range from 2,500 to 80,000 environment steps, so
the arms cannot be compared as they stand; the shortfalls reflect
interrupted runs rather than a design decision. And \textbf{neither arm
reached the solved threshold} within the budget used: the best classical
run finished at 442.7 and the best quantum run at 318.6, against a gate
of 475.

What the runs do establish is that \textbf{the agent learns}: against a
random-policy return of approximately 22 on this environment, the
quantum agent reaches 318.6, so the substituted world model supports
policy improvement through imagination. That is the viability claim, and
it is the one the present data supports.

A properly matched comparison at a sufficient budget is required before
any statement about relative performance, and is listed first in Section
11.

\textbf{{[}RESULTS --- gradient magnitudes at the circuit parameters
over training, from \texttt{prior\_grad\_stats}, demonstrating the
quantum component is optimised rather than carried. The instrumentation
exists; the analysis has not been written.{]}}

\hypertarget{the-world-model-is-used-not-bypassed}{%
\subsubsection{6.3 The world model is used, not
bypassed}\label{the-world-model-is-used-not-bypassed}}

An agent can in principle reach competence while effectively ignoring a
poor world model. We therefore report imagination quality directly:
open-loop prediction error over the imagination horizon, and the
divergence between prior and posterior over training. \textbf{{[}RUN ---
this analysis is not yet implemented and is needed to support the claim
that the quantum prior is doing work.{]}}

\hypertarget{comparison-with-the-classical-prior}{%
\subsubsection{6.4 Comparison with the classical
prior}\label{comparison-with-the-classical-prior}}

\textbf{{[}RUN --- deferred until Section 11.1 is complete. Report final
returns for both arms at a matched budget, with bootstrap intervals, the
number of seeds, and the effect size the sample can resolve.{]}}

\begin{center}\rule{0.5\linewidth}{0.5pt}\end{center}

\hypertarget{design-space}{%
\subsection{7. Design space}\label{design-space}}

\textbf{{[}RUN --- the full ablation grid. See Section 11 for the
planned configurations.{]}}

Planned axes: qubit count, circuit depth, entangler topology (brickwork
against sequential cascade), readout locality, and an
entanglement-removed control in which all two-qubit gates are deleted.
The last is the most important: recent evaluations report that removing
entanglement from quantum models frequently leaves performance unchanged
\textbf{{[}7{]}}, and a first architecture paper should test this of its
own construction rather than leave it to reviewers.

\begin{center}\rule{0.5\linewidth}{0.5pt}\end{center}

\hypertarget{discussion}{%
\subsection{8. Discussion}\label{discussion}}

\textbf{What this establishes.} A parameterised quantum circuit can
serve as the generative dynamics model of a model-based agent, trained
end-to-end through multi-step imagination, with the resulting policy
learning entirely from quantum-generated trajectories. The construction
is stable and reaches competence.

\textbf{What it does not establish.} We make no claim of advantage. At
the scale accessible to exact simulation the quantum prior's function
class is classically reachable, and recent work \textbf{{[}7, 8{]}}
gives good reason to expect parity rather than improvement. Our
comparison is consistent with parity and is not powered to detect small
effects \textbf{{[}RESULTS --- state the detectable effect size{]}}.

\textbf{Why it matters anyway.} The architecture is a prerequisite for
questions that cannot be asked without it: whether quantum dynamics
models help where the dynamics are themselves quantum, whether they
transfer to hardware, and how they behave under finite measurement.
Establishing that the integration is possible, stable and reproducible
is what makes those questions investigable.

\begin{center}\rule{0.5\linewidth}{0.5pt}\end{center}

\hypertarget{limitations}{%
\subsection{9. Limitations}\label{limitations}}

\begin{itemize}
\tightlist
\item
  A single architecture family, at small qubit count, simulated exactly.
  Behaviour under finite shots and device noise is untested.
\item
  \textbf{{[}RESULTS --- number of seeds.{]}} The comparison is reported
  as comparable rather than equal; the sample resolves only large
  effects.
\item
  Evaluation is on low-dimensional control tasks. Scaling to pixel
  observations, where world models are most valuable, is untested.
\item
  The quantum prior carries a substantial wall-clock cost relative to
  the classical control \textbf{{[}RESULTS{]}}.
\item
  \textbf{{[}VERIFY{]}} If the classical control is not
  capacity-matched, the comparison confounds substrate with capacity.
\end{itemize}

\begin{center}\rule{0.5\linewidth}{0.5pt}\end{center}

\hypertarget{reproducibility}{%
\subsection{10. Reproducibility}\label{reproducibility}}

Implementation, verification harness, configurations and run records at
\textbf{{[}AUDIT --- repository URL{]}}. The world model is a faithful
small-scale DreamerV3; the quantum prior is a drop-in substitution, so
the two arms differ in exactly one module.

\begin{center}\rule{0.5\linewidth}{0.5pt}\end{center}

\hypertarget{work-required-before-submission}{%
\subsection{11. Work required before
submission}\label{work-required-before-submission}}

Ordered by importance to the central claim.

\begin{enumerate}
\def\labelenumi{\arabic{enumi}.}
\tightlist
\item
  \textbf{A budget-matched comparison.} The existing runs span 2,500 to
  80,000 environment steps and neither arm met the solved threshold.
  Rerun both arms at a fixed budget, long enough for the classical
  control to solve the task, with at least five seeds. Everything in
  Section 6.4 depends on this, and nothing else on this list substitutes
  for it.
\item
  \textbf{Novelty search.} Systematic, recorded, with search strings and
  dates. The paper's primary claim is a first.
\item
  \textbf{Head-width ablation.} Vary \texttt{prior\_hidden} (64, 16, 8,
  4) to establish whether the circuit or the expansion head carries the
  prediction. Given that the circuit is 0.15\% of the module's
  parameters, this is the ablation a referee will ask for first.
\item
  \textbf{Partial observability.} The \texttt{HideDimensions} wrapper
  exists but has not been run. A history-dependent dynamics model should
  matter most where the observation is incomplete, and this is the
  natural demonstration.
\item
  \textbf{Entanglement-removed control.} Expected by reviewers; cheap.
\item
  \textbf{Ablation grid}: qubits, depth, entangler, readout.
\item
  \textbf{Second environment} (Pendulum) for generality.
\item
  \textbf{Baseline seeds.} The DreamerV3 reference currently has one
  seed.
\item
  \textbf{Imagination-quality analysis} (Section 6.3).
\end{enumerate}

\begin{center}\rule{0.5\linewidth}{0.5pt}\end{center}

\hypertarget{references}{%
\subsection{References}\label{references}}

\textbf{{[}AUDIT --- all citations below require verification of
authors, venue and year.{]}}

{[}1{]} D. Hafner et al.~\emph{Mastering Diverse Domains through World
Models} (DreamerV3).

{[}2{]} S. Y.-C. Chen et al.~\emph{Variational Quantum Circuits for Deep
Reinforcement Learning.}

{[}3{]} \emph{Variational Quantum Circuit Design for Quantum
Reinforcement Learning on Continuous Environments.} arXiv:2312.13798.

{[}4{]} \emph{DreamQAS: Learning a Decision-Useful World Model for
VQE-Efficient Quantum Architecture Search.} arXiv:2607.29491.

{[}5{]} \emph{Quantum Adaptive Self-Attention for Quantum Transformer
Models.} arXiv:2504.05336.

{[}6{]} M. Schuld, R. Sweke, J. J. Meyer. \emph{The effect of data
encoding on the expressive power of variational quantum machine learning
models.} Phys. Rev.~A 103, 032430 (2021).

{[}7{]} J. Bowles, S. Ahmed, M. Schuld. \emph{Better than classical? The
subtle art of benchmarking quantum machine learning models.}
arXiv:2403.07059 (2024).

{[}8{]} V. Liutarevich. \emph{How Many Seeds? Dispersion, Resolving
Power, and a Failed Replication in Quantum Sequence Modelling.} In
preparation. \textbf{{[}Companion paper --- decide whether to cite as
preprint or omit, depending on submission order.{]}}

\end{document}
